% SPDX-License-Identifier: Apache-2.0
% covers: LineFetch
\datasheet{LineFetch}{A host on the bus that reads a run of memory in
bursts, into a channel.}

\dsfacts{\code{lib/parts/src/dma.rs}}{\code{txhdl\_parts::dma::LineFetch}}%
{\code{//docs:examples}}

\subsection*{Function}

A peripheral that wants a lot of memory should not have the core copy
it. The video peripheral's framebuffer is about 420\,000 register
writes for a frame of 720p, and the Ethernet peripheral is the same
shape a byte at a time; in both the processor's whole job during them
is copying.

\code{LineFetch} is the alternative, and it is the half the video and
the MAC both want. Told an address and a count it issues read bursts
and puts the words that come back into a channel. What consumes that
channel is the peripheral's business: for the video it crosses to the
pixel clock through \code{ChanCdc}, and for a MAC it feeds the
transmitter.

\subsection*{Parameters}

\begin{center}\footnotesize
\begin{tabular}{@{}lp{0.68\textwidth}@{}}
\toprule
Parameter & Meaning \\
\midrule
\code{A} & The width of an address, in bits. \\
\code{I} & The width of a burst identifier, in bits. \\
\code{BEATS} & How many beats a burst asks for. \\
\code{WC} & The width of the word count, which says how long a run is. \\
\bottomrule
\end{tabular}
\end{center}

The identifier is a parameter and never a literal, so a link that
widens its identifiers costs this component nothing. The tables below
use \code{LineFetch<32, 2, 16, 16>}, the fetcher of \code{ex\_dma}:
sixteen-beat bursts on a thirty-two bit address.

\dstables{LineFetch}

\subsection*{Behaviour}

\begin{itemize}
\item A run starts when \code{go} is high and nothing is in flight.
\code{base} and \code{words} are read once, there, so a register
written under the engine does not move the run out from under it.
\item A burst goes out when none is in flight and words are still
wanted, asking for \code{BEATS} beats or the rest of the run,
whichever is fewer.
\item A beat is taken when one is offered and the channel out has
room. That room is the backpressure: a consumer that stops taking
stops the bursts.
\item \code{running} is high while a run is unfinished.
\end{itemize}

Two things about the bus are the opposite of what they look like, and
both were written the other way first, and both deadlock.

\emph{The identifier arrives with the burst, not before it.} The host
takes an \code{Issue}, chooses a free identifier and says which on
\code{grant}, so a client sends the burst and catches the identifier
as it comes back. Asking for one and waiting is a deadlock: nothing
is granted until something is issued.

\emph{A read burst is finished by its last beat, not by \code{done}.}
The host puts a write's response on \code{done} and answers a read on
the beat channel. So the identifier goes back as the last beat is
taken, and the last beat is only taken when there is room to give it
back.

\subsection*{Verification}

\begin{itemize}
\item \code{ex\_dma} fetches sixty-four words through sixteen-beat
bursts across a real link, with \code{AxiHost}, \code{AxiPer} and a
memory on the far side. The build checks the netlist against that run
under nvc and Verilator:
\code{//docs:dma\_sim\_linefetch\_tb\_test} and
\code{//docs:dma\_vsim\_test}.
\item The example asserts rather than prints. Every word arrived;
each word is the one at its own address, since the memory holds its
address in each word, so a word from the wrong burst or out of order
says so; and the run came in under a cycle bound a path spending an
address phase per word could not meet. Sixty-four words in 93 cycles
against a bound of 128.
\item The memory model there answers bursts rather than words. A
model that answered one beat per request would make a burst look like
a word: the words would still all arrive, the cycle count would still
pass, and nothing about bursting would be demonstrated.
\end{itemize}

\subsection*{Choosing \code{BEATS}}

\code{BEATS} is the lever that decides whether real memory latency
matters, and the number that looks like a measurement above does not
say so on its own.

The 93 cycles are measured against a memory model that answers a beat
the cycle after the request: no activation, no CAS, no refresh. So
that figure is a floor and not an estimate. What a real memory adds is
one first-beat latency per burst, and the burst count is
\code{ceil(words / BEATS)}, so the exposed latency is the burst count
multiplied by whatever the controller's latency is. One burst is in
flight at a time, so none of it is hidden.

That makes the parameter, rather than the engine, the answer to a
consumer that is faster than a raster. A frame of 1500 bytes is 375
words, and at gigabit the wire takes about 1200 cycles at 100\,MHz:

\begin{center}\footnotesize
\begin{tabular}{@{}lll@{}}
\toprule
\code{BEATS} & Bursts for a frame & Fits 1200 cycles while \\
\midrule
16 & 24 & latency is under about 27 cycles \\
256 & 2 & latency is under about 327 cycles \\
\bottomrule
\end{tabular}
\end{center}

A realistic DDR3 first-beat latency is tens of cycles, so sixteen-beat
bursts can lose that margin and 256-beat bursts cannot. The raster
uses sixteen because its consumer takes a pixel per visible cycle and
a longer burst buys it nothing; a MAC should use a longer one. Both
are the same component with a different parameter, and neither needs a
second burst in flight, which is what keeps a second identifier, a
tracker and out-of-order answers out of the design.

\subsection*{Limits}

\textbf{It counts the beats of a burst itself and never reads the
returned beat's \code{last}.} So it cannot notice a peripheral that
answers fewer beats than the burst asked for: it takes what arrives
and waits for the rest, holding \code{running} high, with nothing to
say why. Against the Wishbone bridge before issue 471, which answered
any read with one beat, sixty four words became one and the run did
not end.

That is worth knowing before joining it to a peripheral written
elsewhere. The engine assumes AXI4 is obeyed, and a peripheral that
under-delivers is outside what it can detect rather than something it
reports.


\textbf{A burst must not cross a 4\,KB address boundary}, which is
AXI4's rule and not this component's, and nothing here checks it.
\code{BEATS} of 256 is 1\,KB, so a run whose base is 1\,KB aligned
satisfies it without a special case: every burst but the last is a
full 1\,KB from a 1\,KB aligned address, and the last is shorter from
one. A caller that starts a run at an arbitrary address with a long
\code{BEATS} can violate the rule, and the symptom is a burst the
interconnect is entitled to answer however it likes.

One burst is in flight at a time, deliberately. A second would want a
second identifier, a tracker to tell the two apart and an answer for
what happens when they come back out of order, and none of that buys
anything while the consumer is a raster taking one pixel per visible
cycle.

It reads and never writes. \code{done} is drained so that a write
answer arriving from elsewhere cannot block the host, and is
otherwise unused.

It runs on one clock, \code{DefaultClock}, which is the bus clock.
Crossing to a consumer on another clock is \code{ChanCdc}'s job and
not this component's.

Beats are words: \code{size} is fixed at two, four bytes a beat.
